\section{Results}

The benchmark instances used in the experiments were the provided LOP instances of sizes 150 and 250. Each algorithm was applied once to each instance, and a random seed was set for reproducibility and fair comparaison.
For each run, two values were recorded:
\begin{itemize}[leftmargin=1.5em]
    \item the final solution cost
    \item the computation time
\end{itemize}

Below is a summary table, where solution quality is measured using the relative percentage deviation from the best known solution, defined as
\[
\Delta_{ki} = 100 \cdot \frac{\text{best-known}_i - \text{cost}_{ki}}{\text{best-known}_i}.
\]

With this convention, smaller values correspond to better solutions.

\subsection{Exercise 1.1 : Summary of the 12 Iterative Improvement Algorithms}
\begin{table}[H]
\centering
\caption{Summary results for the 12 iterative improvement algorithms.}
\label{tab:ex1}
\begin{tabular}{lccc}
\toprule
\textbf{Algorithm Variant} & \textbf{Avg. Deviation (\%)} & \textbf{Std. Dev.} & \textbf{Total Time (s)} \\
\midrule
cw best exchange      & 3.8465  & 0.6054 & 32.9845 \\
cw best insert        & 2.2608  & 0.4065 & 33.7210 \\
cw best transpose     & 17.0783 & 1.6065 & 0.2194 \\
cw first exchange     & 2.9660  & 0.4613 & 186.7398 \\
cw first insert       & 2.0164  & 0.4109 & 128.2090 \\
cw first transpose    & 17.0863 & 1.6042 & 0.2148 \\
random best exchange  & 3.6027  & 0.5048 & 35.8138 \\
random best insert    & 2.3037  & 0.4635 & 37.9277 \\
random best transpose & 34.4701 & 3.8033 & 0.0035 \\
random first exchange & 2.8159  & 0.4810 & 228.0147 \\
random first insert   & 2.0211  & 0.4507 & 189.8870 \\
random first transpose& 34.6058 & 3.7984 & 0.0024 \\
\bottomrule
\end{tabular}
\end{table}

\subsection{Exercise 1.2 : Summary of the VND Algorithms}

\begin{table}[H]
\centering
\caption{Summary results for the two VND algorithms.}
\label{tab:ex2}
\begin{tabular}{lccc}
\toprule
\textbf{Method} & \textbf{Avg. Deviation (\%)} & \textbf{Std. Dev.} & \textbf{Total Time (s)} \\
\midrule
vnd-tei & -2.0633 & 0.4191 & 109.6544 \\
vnd-tie & -1.9959 & 0.4220 & 80.5438 \\
\bottomrule
\end{tabular}
\end{table}

\subsection{Tests}

A series of tests were performed to compare the algorithms, to answer the questions of Is there a statistically significant difference between the solution quality generated by the different
algorithms? 

The wilcoxon 




The two VND variants were compared using paired Student t-tests and Wilcoxon signed-rank tests.

\begin{itemize}[leftmargin=1.5em]
    \item Average deviation vnd-tei: \(-2.0633\)
    \item Average deviation vnd-tie: \(-1.9959\)
    \item Std deviation vnd-tei: \(0.4191\)
    \item Std deviation vnd-tie: \(0.4220\)
    \item Paired t-test p-value: \(0.2136\)
    \item Wilcoxon p-value: \(0.4169\)
\end{itemize}

Although the transpose--insert--exchange variant obtained a slightly better average deviation than the transpose--exchange--insert variant, both statistical tests returned p-values greater than \(0.05\). Therefore, no statistically significant difference in solution quality was observed between the two VND algorithms.